
\begin{center}
\LARGE{\textbf{Défi de Conversions Numériques : Le Déchiffrement de Données}}
\end{center}


Vous êtes un ingénieur en informatique travaillant sur un ancien système qui utilise trois bases numériques différentes pour stocker et traiter l'information. Pour que le système fonctionne correctement, vous devez vous assurer que toutes les données sont interprétées correctement.

\textbf{Votre mission :} Convertir les données brutes fournies dans le format requis.

\vspace{0.3cm}
\hrule
\vspace{0.3cm}

\section*{Partie I : Le Déboguage des Registres Binaires}

Les registres CPU affichent des valeurs en binaire. Vous devez savoir quelle valeur décimale et hexadécimale ils représentent.

\begin{center}
\begin{tabular}{|c|c|c|c|}
    \hline
    \textbf{Question} & \textbf{Registre Binaire} & \textbf{Convertir en Décimal } & \textbf{Convertir en Hexadécimal } \\
    \hline
    \textbf{1} & $10101010_2$ & & \\
    \hline
    \textbf{2} & $11100101_2$ & & \\
    \hline
    \textbf{3} & $11010000_2$ & & \\
    \hline
\end{tabular}
\end{center}

\vspace{0.5cm}

\section*{Partie II : Les Adresses Mémoire Décimales}

Vous devez paramétrer des sauts dans la mémoire en utilisant des adresses hexadécimales, mais les spécifications sont fournies en décimal.

\begin{center}
\begin{tabular}{|c|c|c|c|}
    \hline
    \textbf{Question} & \textbf{Adresse Décimale } & \textbf{Convertir en Binaire} & \textbf{Convertir en Hexadécimal} \\
    \hline
    \textbf{4} & $120_{10}$ & & \\
    \hline
    \textbf{5} & $255_{10}$ & & \\
    \hline
    \textbf{6} & $16_{10}$ & & \\
    \hline
\end{tabular}
\end{center}

\vspace{0.5cm}

\section*{Partie III : Les Codes d'Opération Hexadécimaux}

Les codes d'opération machine (\textit{opcodes}) sont souvent représentés en hexadécimal. Pour vérifier leur validité, vous devez les convertir.

\begin{center}
\begin{tabular}{|c|c|c|c|}
    \hline
    \textbf{Question} & \textbf{Code Hexadécimal} & \textbf{Convertir en Décimal } & \textbf{Convertir en Binaire } \\
    \hline
    \textbf{7} & $A5_{16}$ & & \\
    \hline
    \textbf{8} & $3F_{16}$ & & \\
    \hline
    \textbf{9} & $C0_{16}$ & & \\
    \hline
\end{tabular}
\end{center}

\vspace{0.5cm}

\section*{Question Bonus : L'Erreur de Parité}

\textbf{10.} Un paquet de données a été reçu avec l'octet binaire suivant : $11010011_2$. Quelle est sa valeur \textbf{décimale} ? Et quelle est sa valeur \textbf{hexadécimale} ?


% \section*{Corrigé du Défi}

% \begin{center}
% \begin{tabular}{|c|c|c|c|c|}
%     \hline
%     \textbf{Question} & \textbf{Base Initiale} & \textbf{Binaire ($\text{Base 2}$)} & \textbf{Décimal ($\text{Base 10}$)} & \textbf{Hexadécimal ($\text{Base 16}$)} \\
%     \hline
%     \textbf{1} & $10101010_2$ & $10101010_2$ & $170_{10}$ & $AA_{16}$ \\
%     \hline
%     \textbf{2} & $11100101_2$ & $11100101_2$ & $229_{10}$ & $E5_{16}$ \\
%     \hline
%     \textbf{3} & $11010000_2$ & $11010000_2$ & $208_{10}$ & $D0_{16}$ \\
%     \hline
%     \textbf{4} & $120_{10}$ & $01111000_2$ & $120_{10}$ & $78_{16}$ \\
%     \hline
%     \textbf{5} & $255_{10}$ & $11111111_2$ & $255_{10}$ & $FF_{16}$ \\
%     \hline
%     \textbf{6} & $16_{10}$ & $00010000_2$ & $16_{10}$ & $10_{16}$ \\